\documentclass[11pt]{article}

% ACL style (review submission)
\usepackage[review]{acl}

% Standard package includes
\usepackage{times}
\usepackage{latexsym}

\usepackage[T1]{fontenc}
\usepackage[utf8]{inputenc}
\usepackage{microtype}
\usepackage{inconsolata}
\usepackage{graphicx}

% Additional packages for this paper
\usepackage{hyperref}
\usepackage{url}
\usepackage{booktabs}
\usepackage{amsfonts}
\usepackage{amsmath}
\usepackage{amssymb}
\usepackage{nicefrac}
\usepackage{xcolor}
\usepackage{multirow}
\usepackage{subcaption}
\usepackage{algorithm}
\usepackage{algorithmic}
\usepackage{tablefootnote}
\usepackage{pifont}
\newcommand{\cmark}{\ding{51}}
\newcommand{\xmark}{\ding{55}}

\setlength\titlebox{6cm}

\title{Supplementary Material: Structured Social State for Auditable NPC Dialogue Generation}

\author{Anonymous \\ Affiliation / Address line 1 \\ Affiliation / Address line 2 \\ Affiliation / Address line 3 \\ \texttt{email@domain} \\And Anonymous \\ Affiliation / Address line 1 \\ Affiliation / Address line 2 \\ Affiliation / Address line 3 \\ \texttt{email@domain} \\}

\begin{document}

\maketitle

\appendix

\section{Real-World Use Cases and Practical Accomplishments}
\label{app:usecases}

This appendix addresses the natural question: \textit{what can someone actually do with these models?}
We describe concrete deployment scenarios for each track and the practical capabilities the trained models provide.

\subsection{Track A: From-Scratch SLMs}

\textbf{What they are:} Six small language models (15--45M parameters) trained entirely from random initialization on 16,905 lines of NPC dialogue.
\textbf{What they accomplish:} They learn to generate domain-appropriate NPC speech---guards sound like guards, merchants sound like merchants---without requiring any pretrained model, cloud API, or internet connection.

\textbf{Real-world use cases:}
\begin{enumerate}
    \item \textbf{Embedded game engines.} A 16M-parameter model is approximately 60MB on disk and runs at interactive speeds on a CPU.
    This means you can ship NPC dialogue generation inside a game binary with no external dependencies.
    No GPU, no network, no API costs---just a compiled model running locally.
    \item \textbf{Privacy-preserving applications.} Because these models are trained from scratch on in-domain data, they never see internet text, user conversations, or any external data.
    For healthcare training simulations or corporate role-play scenarios where data must never leave the device, from-scratch SLMs provide a fully offline, inspectable solution.
    \item \textbf{Architecture research platform.} The 15--16M parameter scale makes these models cheap to train (20 epochs in under 5 minutes on an A30) and easy to modify.
    Researchers can test new architectures (new attention mechanisms, new SSM variants, new expert-routing strategies) within hours rather than weeks; statistical claims still require multi-seed reruns.
\end{enumerate}

\textbf{Current capability:} The MoE model achieves mean val\_ppl $=$ 43.38 across three seeds (best seed 42.39) on held-out dialogue.
This is not yet human-quality, but it demonstrates that domain-specific, small-scale models can capture NPC speaking patterns without pretraining.
\textbf{Inference cost:} CPU-only, $\approx$60\,MB on disk, no GPU required.
With more data and a few architectural improvements, this track could produce production-quality embedded dialogue models.

\subsection{Track B: Personality and Affect Encoders}

\textbf{What they are:} Two DistilBERT-based regression models that predict personality traits (OCEAN) from NPC profile text and emotional state (VAD) from dialogue context.
\textbf{What they accomplish:} They provide a computational \textit{theory of mind} for NPCs---a machine-readable estimate of who a character is and how they feel at any moment.

\textbf{Real-world use cases:}
\begin{enumerate}
    \item \textbf{Automated NPC authoring.} Game writers define NPCs with short text descriptions (``a paranoid guard who values duty over mercy'').
    The personality encoder converts this to a 5-dimensional OCEAN vector that downstream systems can use to modulate dialogue style, decision-making, and animation parameters (e.g., extraverted NPCs gesticulate more; neurotic NPCs have wider eye movement).
    \item \textbf{Live social perception.} As a conversation unfolds, the affect encoder tracks the NPC's emotional trajectory in real time.
    A sudden spike in arousal or drop in valence can trigger gameplay events: the NPC might become aggressive, flee, or call for backup.
    \item \textbf{Player modelling.} The same encoder architecture, applied to player dialogue, could estimate the player's emotional state and adjust game difficulty, narrative branching, or NPC behaviour accordingly.
\end{enumerate}

\textbf{Current capability:} Personality encoder achieves val\_f1 $=$ 0.678; affect encoder achieves val\_ccc $=$ 0.559.
A 414-NPC personality cache is precomputed and available for instant lookup.
The encoders run in under 10ms per inference on CPU.

\subsection{Track C: Response Generation Models}

\textbf{What they are:} Three fine-tuned language models that generate NPC dialogue responses, ranging from a conditioned 1.1B-parameter model (ConditionalDialogue) to a 5.1B-parameter mixture-of-experts (Gemma 4 E2B).
\textbf{What they accomplish:} They produce fluent, character-consistent NPC responses that respect the NPC's personality and emotional state.

\textbf{Real-world use cases:}
\begin{enumerate}
    \item \textbf{Dynamic NPC dialogue in games.} Instead of pre-writing every possible NPC response, a game can generate responses on the fly conditioned on who the NPC is (OCEAN cache) and how they currently feel (VAD encoder).
    This enables NPCs that react differently to the same player query depending on context---a capability that hand-authored dialogue trees struggle to provide at scale.
    \item \textbf{Social-state-aware chatbots.} Beyond games, any conversational agent that needs to maintain a consistent persona can benefit from explicit personality and affect conditioning.
    Customer service bots, therapeutic agents, and educational tutors all have personas that should remain stable across interactions.
    \item \textbf{Baseline for stronger models.} TinyLlama SFT (val\_ppl $=$ 3.30) provides a clean baseline: ``this is how well a pretrained 1.1B model does without any social state.''
    The 12.3\% improvement from adding conditioning quantifies the value of the prefix mechanism, while placebo controls show that OCEAN/VAD values should not be treated as the causal signal.
    \item \textbf{Gemma as an engineering baseline.} The Gemma-2-2B-it run provides the corrected scientific baseline in this draft, while the Gemma 4 E2B run demonstrates that sparse MoE architectures can be fine-tuned on a single 24\,GB GPU for training.
    Inference requires 12--16\,GB VRAM with 4-bit quantisation---well within reach of a consumer RTX 4080 or cloud T4 instance.
    For applications that can afford a larger inference budget, Gemma-family models remain useful baselines, but the current Gemma runs do not outperform the task-specific conditional model on this validation split.
\end{enumerate}

\textbf{Current capability:} ConditionalDialogue achieves mean val\_ppl $=$ 2.88 across three seeds, with zero gated leakage in the Qwen3-4B pipeline.
The 12.3\% gain over the unconditioned baseline should be interpreted as a soft-prefix capacity effect because placebo controls reach the same PPL range.

\subsection{Track D: Structured LLM Pipeline}

\textbf{What it is:} A three-stage Qwen3-4B pipeline that (1) predicts the full 29-dimension social state from dialogue context, (2) generates responses, and (3) jointly trains both objectives with a consistency loss.
\textbf{What it accomplishes:} It provides an end-to-end system where a single model can both \textit{understand} the social dynamics of a conversation and \textit{act} on that understanding.

\textbf{Real-world use cases:}
\begin{enumerate}
    \item \textbf{Socially-aware NPC AI prototype.} This is the closest component in the project to an integrated NPC social-reasoning loop.
    Given conversation history, the latent predictor estimates the NPC's full social state---what they think of the player, whether they are under pressure to keep secrets, what their strategic intent is.
    The response generator then produces dialogue that is consistent with that state.
    If the state says ``secrecy\_pressure $=$ high, reveal\_decision $=$ none,'' the system can route or constrain generation through an inspectable disclosure policy.
    \item \textbf{Game master AI.} The 29-head predictor could serve as a ``social radar'' for a game master AI that tracks the social state of every NPC in a scene and decides when to trigger narrative events, introduce new characters, or escalate conflicts.
    \item \textbf{Explainable NPC behaviour.} Because the latent state is a discrete, human-readable vector (not a hidden embedding), game designers can inspect why an NPC behaved a certain way.
    ``The NPC refused to answer because secrecy\_pressure was high'' is a debuggable explanation that an embedding vector cannot provide.
    \item \textbf{Fast/slow routing in production.} The rule-based router (Section~5.3 of the main paper) classifies each turn as ``fast path'' (simple responses, no social complexity) or ``slow path'' (requires full $Z_t$ inference).
    In a game with hundreds of NPCs, most interactions are fast-path (a shopkeeper saying ``That'll be 10 gold''), reserving the expensive latent predictor for socially complex encounters.
\end{enumerate}

\textbf{Current capability:} Latent predictor (Qwen3-4B) achieves response\_policy\_f1 $=$ 0.622, reveal\_decision\_f1 $=$ 0.656, and mean accuracy $=$ 0.688 (true $\kappa{=}$0.484).
Response generator achieves ROUGE-L $=$ 0.145 with length ratio $=$ 0.91, distinct-2 $=$ 0.638, and zero gated secret leakage under the heuristic detector.
Joint model trains end-to-end (train loss $=$ 1.95, val loss $=$ 3.89 with consistency weight $\lambda{=}0.5$).
All three stages have been trained and independently evaluated; the joint model will be released upon acceptance.

\subsection{What We Accomplish Overall}

Table~\ref{tab:accomplish} (below) summarises the concrete capabilities delivered by this work.

\begin{table}[h]
\centering
\caption{Concrete accomplishments of the trained models.}
\label{tab:accomplish}
\small
\begin{tabular}{@{}p{0.18\linewidth}p{0.35\linewidth}p{0.35\linewidth}@{}}
\toprule
\textbf{Capability} & \textbf{How We Demonstrate It} & \textbf{Who Benefits} \\
\midrule
Offline NPC dialogue (no pretrained model, no API, 60MB on disk)
    & Track A: 4 from-scratch SLMs trained on 16,905 lines; MoE mean PPL $=$ 43.38 (3 seeds)
    & Game developers evaluating embedded dialogue; privacy-critical applications \\
\midrule
Computational theory of mind (personality $+$ affect from text)
    & Track B: OCEAN encoder (f1 $=$ 0.678) $+$ VAD encoder (ccc $=$ 0.559); 414-NPC cache
    & Game writers doing automated NPC authoring; live social perception systems \\
\midrule
Soft-prefix dialogue with 12.3\% PPL improvement
    & Track C: ConditionalDialogue (mean ppl $=$ 2.88, 3 seeds) vs. TinyLlama SFT (ppl $=$ 3.30); placebo controls show the gain is prefix-capacity rather than OCEAN/VAD semantics
    & Any conversational agent needing consistent persona and information control \\
\midrule
End-to-end social perception $+$ generation in one model
    & Track D: 3-stage Qwen3 pipeline; joint model trains both objectives simultaneously
    & Production NPC AI; game master systems; explainable AI for games \\
\midrule
Practical deployability across hardware tiers
    & Tracks A/B: CPU-only inference (60\,MB on disk, {$<$}10\,ms). Tracks C/D: 4--16\,GB VRAM inference with 4-bit quantisation. Training: 8--24\,GB depending on model.
    & Independent researchers, small studios; incrementally deployable from CPU to single consumer GPU \\
\bottomrule
\end{tabular}
\end{table}

In summary, this work delivers not just a set of benchmark numbers, but a practical toolkit spanning the full hardware spectrum: CPU-only models for embedded deployment (Track~A, {$<$}60\,MB), lightweight encoders for social perception (Track~B, {$<$}10\,ms on CPU), conditioned generators for character-consistent dialogue (Track~C, 4--16\,GB inference), and a full pipeline for integrated social-state-aware NPC AI (Track~D).
Each component can be used independently or composed into larger systems.
Training requires at most a single 24\,GB GPU; inference requirements are substantially lower, with Tracks~A and~B requiring no GPU at all.

\section{System Data Flow}
\label{app:dataflow}

This appendix documents the input $\rightarrow$ model $\rightarrow$ output flow for each architectural component in the benchmark.
All dimensions and data paths correspond to the actual training configuration used in the main paper.

\subsection{Track A: From-Scratch Small Language Models}

\paragraph{GPT, MoE, Mamba-like (unconditioned).}
\textbf{Input:} Plain-text dialogue from \texttt{data/dialogue/train.txt} (16,905 lines, 414 unique NPCs).
Each line is tokenised with the GPT-2 tokeniser (vocabulary size 50,257), truncated or padded to sequence length 256.
\textbf{Model:} Architecture-specific backbone (GPT: 6-layer transformer decoder, MoE: 4-experts with top-2 routing, Mamba-like: selective-scan SSM), embedding dimension 512, 8 attention heads where applicable.
\textbf{Output:} Next-token logits over vocabulary; training loss is cross-entropy.
Evaluation metric: validation perplexity (Table~3 in the main paper).

\paragraph{PrefixGPT (conditioned).}
\textbf{Input:} Same dialogue text as above, plus an 8-dimensional conditioning vector $c_t = [\text{OCEAN}_5, \text{VAD}_3]$ concatenated from the personality and affect encoders.
\textbf{Model:} 6-layer GPT with a soft-prefix module: a 2-layer MLP projects $c_t \in \mathbb{R}^8$ to $\ell \cdot d$ prefix embeddings ($\ell=4$ prefix tokens, $d=512$ hidden), which are prepended to the input sequence.
\textbf{Output:} Next-token logits; val\_ppl $=$ 44.11 (mean, 3 seeds).

\subsection{Track B: Conditioning Encoders}

\paragraph{Personality Encoder (OCEAN).}
\textbf{Input:} NPC profile text from \texttt{data/npc\_profiles.csv} (414 entries, e.g., ``A formal guard who values duty and secrecy'').
Tokenised with \texttt{distilbert-base-uncased}, max length 512.
\textbf{Model:} DistilBERT (66M) $\rightarrow$ mean-pooled hidden state $\rightarrow$ Linear(768, 5).
MSE loss over 5 OCEAN dimensions.
Best epoch 4/15.
\textbf{Output:} 5-dim OCEAN vector per NPC, cached to \texttt{artifacts/personality\_cache.jsonl}.
Validation: val\_f1 $=$ 0.678, val\_mse $=$ 0.248.

\paragraph{Affect Encoder (VAD).}
\textbf{Input:} Dialogue context text from \texttt{data/affect/train.csv} (500 synthetic samples).
\textbf{Model:} DistilBERT (66M) $\rightarrow$ mean-pooled hidden state $\rightarrow$ Linear(768, 3).
Combined MSE $+$ CCC loss over 3 VAD dimensions.
Best epoch 13/15.
\textbf{Output:} 3-dim VAD vector per dialogue turn.
Validation: val\_ccc $=$ 0.559, val\_mse $=$ 0.005.

\subsection{Track C: Response Generation}

\paragraph{ConditionalDialogue (OCEAN + VAD conditioned).}
\textbf{Input:} Dialogue history $+$ OCEAN vector (from cache) $+$ VAD vector (live from affect encoder) $\rightarrow$ 8-dim $c_t$.
\textbf{Model:} TinyLlama-1.1B with LoRA ($r=8$, $\alpha=16$) $+$ ConditionalSoftPrefix: $c_t$ projected to 4 soft-prefix tokens via 2-layer MLP, prepended to input embeddings.
The prefix encoder and LoRA adapter are trained jointly.
\textbf{Output:} Generated NPC response.
Best val\_ppl $=$ 2.88 (mean across seeds 42--44) at epoch 5/5.

\paragraph{TinyLlama SFT (baseline).}
\textbf{Input:} Dialogue history only (no personality or affect conditioning).
\textbf{Model:} TinyLlama-1.1B with LoRA ($r=16$, $\alpha=32$). Standard supervised fine-tuning: input $\rightarrow$ target response, with input tokens masked (label $=-100$).
\textbf{Output:} Generated NPC response.
Best val\_ppl $=$ 3.30 at epoch 1/3.
Conditioning gain $=$ 12.3\%.

\paragraph{Gemma 4 E2B + QLoRA.}
\textbf{Input:} Structured dialogue records from \texttt{from\_gen\_train.jsonl}, each containing \texttt{npc\_id}, \texttt{npc\_profile}, \texttt{dialogue\_context}, and \texttt{target\_response}.
Formatted via the Gemma chat template: \texttt{<start\_of\_turn>user\textbackslash n[NPC profile]\textbackslash n\textbackslash n[player query]<end\_of\_turn>}.
\textbf{Model:} Gemma 4 E2B (16B total, 2B active MoE) with 4-bit nf4 QLoRA, LoRA $r=16$, $\alpha=16$, 37.9M trainable parameters (0.74\% of total).
Uses \texttt{target\_modules="all-linear"} to handle \texttt{Gemma4ClippableLinear} layers.
Memory: \texttt{PYTORCH\_CUDA\_ALLOC\_CONF=expandable\_segments:True} needed for evaluation on 24\,GB GPUs; training fits without tuning.
Inference: 12--16\,GB VRAM with 4-bit quantisation.
Training via \texttt{SFTTrainer} with \texttt{formatting\_func} and \texttt{processing\_class} (trl $\geq$ 1.3.0 API).
\textbf{Output:} Generated NPC response.
Best val\_ppl $=$ 16.24 at epoch 1/2 (exploratory 1-epoch run; rereun at 3 epochs achieves 13.48, see Table~4 in the main paper).

\subsection{Track D: Structured LLM Pipeline}

\paragraph{Stage 1: Latent State Predictor.}
\textbf{Input:} Dialogue context from \texttt{data/splits/train.jsonl} with full 29-field $Z_t$ labels.
Context includes scene description and conversation history, tokenised to max 1,024 tokens.
\textbf{Model:} Qwen3-4B backbone $+$ LoRA $+$ 29 parallel classification heads.
Each head is a 2-layer MLP (hidden $\rightarrow$ 256 $\rightarrow$ $n_\text{classes}$) applied to the last-token hidden state.
Weighted multi-head loss: $\lambda_R = \lambda_D = 2.0$, $\lambda_M = 1.5$, others $=1.0$.
\textbf{Output:} 29 discrete predictions per turn.
Best response\_policy\_f1 $=$ 0.622, reveal\_decision\_f1 $=$ 0.656, mean\_accuracy $=$ 0.688 at epoch 5/5.

\paragraph{Stage 2: Response Generator (SFT).}
\textbf{Input:} \texttt{data/splits/train\_sft.jsonl} with \texttt{input} (scene $+$ history) and \texttt{target} (NPC response).
Tokenised as \texttt{input $+$ "\textbackslash n" $+$ target}, max length 2,048.
Input portion masked (label $=-100$).
\textbf{Model:} Qwen3-4B $+$ LoRA ($r=32$, $\alpha=64$). Standard causal LM loss on target tokens only.
\textbf{Output:} Generated NPC response.
Best val\_loss $=$ 0.044 at epoch 3/3.

\paragraph{Stage 3: Joint Model.}
\textbf{Input:} Both data sources combined: each turn has dialogue context, target response, and 29 $Z_t$ labels.
\textbf{Model:} Shared Qwen3-4B backbone with both the 29-head predictor and the LM head active simultaneously.
Joint loss: $\mathcal{L}_\text{joint} = \mathcal{L}_\text{heads} + \mathcal{L}_\text{LM}$.
Gradients from both objectives flow through shared parameters.
\textbf{Output:} Both social state predictions and generated responses from a single model.
Train loss 1.95, val loss 3.89 at epoch 3/3.

\subsection{End-to-End Integration}

Figure~2 in the main paper shows the complete flow from scenario generation through training to evaluation.
The conditioning path (Track B $\rightarrow$ Track C) operates as: NPC profile $\rightarrow$ Personality encoder $\rightarrow$ OCEAN cache $\rightarrow$ concatenated with live VAD from Affect encoder $\rightarrow$ 8-dim $c_t$ $\rightarrow$ ConditionalSoftPrefix $\rightarrow$ TinyLlama with LoRA $\rightarrow$ response.
The structured LLM path (Track D) operates as: dialogue history $\rightarrow$ Qwen3-4B latent predictor $\rightarrow$ 29-dim $\hat{Z}_t$ $\rightarrow$ Qwen3-4B response generator $\rightarrow$ response.
In the joint model, both prediction and generation share a single backbone forward pass, producing $\hat{Z}_t$ and the response simultaneously.

\section{Training Pipeline Infrastructure \& Code Modifications}
\label{app:infrastructure}

Training 12 models across four tracks with three deep learning frameworks (transformers, trl, peft) and two separate Python virtual environments required substantial infrastructure development and debugging.
We document the key modifications here for reproducibility.

\subsection{Slurm Job Management}

All training ran on an HPC cluster with NVIDIA A30 (24GB) and NVIDIA L20 (48GB) nodes managed by Slurm.
We developed eight Slurm scripts:

\begin{itemize}
    \item \texttt{slurm\_train.sh}: Single-model training for any SLM or LLM stage. Accepts \texttt{\{slm,llm\}} and stage parameters. Includes environment activation, CUDA module loading, and working directory management.
    \item \texttt{slurm\_array.sh}: Job arrays for multi-architecture, multi-seed hyperparameter sweeps.
    \item \texttt{slurm\_eval.sh}: Standalone evaluation with metrics bundle generation.
    \item \texttt{slurm\_train\_eval.sh}: Combined training $+$ auto-evaluation in a single job.
    \item \texttt{resume\_training.sh}: Full pipeline orchestrator with pre-flight verification (1-epoch GPT test before full training).
    \item \texttt{submit\_test\_jobs.sh}: Submits 12 lightweight test jobs (1 epoch, 15 min each) covering every training task.
    \item \texttt{submit\_missing.sh}: Auto-detects incomplete runs via checkpoint inspection and submits only pending jobs.
    \item \texttt{gpu\_test.sh}: Interactive GPU session for environment verification.
\end{itemize}

Cluster-specific conventions: \texttt{--gpus-per-node=N} (not \texttt{--gres=gpu}), \texttt{--ntasks-per-node=1}.
The codebase uses symlinks to fast NVMe scratch storage for I/O on compute nodes.

\subsection{LLM Training Configuration (Qwen3)}

The three-stage Qwen3 pipeline uses the following config structure (\texttt{llm\_finetuning/configs/}):

\begin{itemize}
    \item \texttt{train\_latent.yaml}: Qwen3-4B. 4-bit QLoRA, \texttt{nf4} quantisation, \texttt{bfloat16} compute dtype. 29 classification heads. Focal loss enabled via \texttt{focal\_gamma} for weak heads.
    \item \texttt{train\_response.yaml}: Qwen3-4B SFT on dialogue data. \texttt{max\_seq\_len: 2048} (increased from 1024 to prevent NaN loss from truncated labels). LoRA $r=32$, $\alpha=64$. Secret-span masking uses per-record \texttt{secret\_strings} from packaged SFT split.
    \item \texttt{train\_joint.yaml}: Joint training combining latent prediction and response generation losses.
    \item \texttt{eval.yaml}: Length-aware response generation, dual leakage accounting, and gold/predicted routing modes. Predicted routing requires \texttt{predicted\_zt.jsonl} from latent evaluation.
\end{itemize}

\subsection{Evaluation Framework}

The intended comprehensive evaluation script computes the following metrics when generated predictions and labels are available:

\begin{itemize}
    \item \textbf{Per-head metrics:} Accuracy, true Cohen's $\kappa$, macro-F1, MCC for all 28 single-label heads.
    \item \textbf{Per-group aggregation:} Mean $\pm$ standard deviation across $C_t$, $A_t$, $M_t$, $R_t$, $N_t$, $D_t$ groups.
    \item \textbf{Response quality:} ROUGE-L with bootstrap 95\% CI ($n=1{,}000$ resamples), BLEU-1/2/4, Distinct-1/2/3, 3-gram and 5-gram repetition rates, generation-to-reference length ratio.
    \item \textbf{Policy consistency:} Keyword-based checks for withhold, reveal, and soothe adherence.
    \item \textbf{Secret leakage:} Regular expression matching against known secret patterns from scenario definitions, reported both as gated leakage over \texttt{reveal\_decision=none} turns and ungated leakage across all sampled turns.
    \item \textbf{Contradiction detection:} Checks for contradictory statements within single generated responses.
    \item \textbf{Routing:} Precision, recall, F1, slow-path rate, and prediction coverage for the fast/slow path router. Gold-state routing is a deterministic sanity check; predicted-state routing consumes \texttt{predicted\_zt.jsonl} and is the deployment-relevant mode.
\end{itemize}

Evaluation artifacts (\texttt{comprehensive\_results.json}, \texttt{per\_head\_metrics.csv}, \texttt{paper\_tables.md}) are included in the repository.

\subsection{Code Modifications for Reproducibility}

Table~\ref{tab:codefixes} (below) documents the 16 code modifications required to train all 15 models.
These represent real infrastructure challenges that future researchers working with similar multi-framework, multi-model pipelines will likely encounter.

\begin{table*}[t]
\centering
\caption{Code modifications required for the complete training pipeline.}
\label{tab:codefixes}
\scriptsize
\begin{tabular}{@{}p{0.28\linewidth}p{0.70\linewidth}@{}}
\toprule
\textbf{Issue} & \textbf{Resolution} \\
\midrule
\texttt{build\_personality\_cache} missing in \texttt{build\_caches.py} & Used \texttt{encode\_profiles} in \texttt{smoke\_test.py} \\
\midrule
Tokenizer path returned \texttt{NoneType} on affect encoder load & Hardcoded \texttt{distilbert-base-uncased} tokenizer \\
\midrule
\texttt{torch.cuda.amp.GradScaler} deprecated (PyTorch 2.6) & Migrated to \texttt{torch.amp.GradScaler('cuda', ...)} \\
\midrule
\texttt{get\_sentence\_embedding\_dimension} renamed in \texttt{sentence-transformers} & Added fallback to \texttt{get\_embedding\_dimension} in \texttt{memory\_store.py} \\
\midrule
\texttt{faiss} not installed in SLM virtual environment & \texttt{pip install faiss-cpu} \\
\midrule
SLM training looked for data in wrong directory & Added \texttt{cd slm\_training/} in \texttt{slurm\_train.sh} \\
\midrule
Condition vectors (float32) vs. embeddings (BFloat16) dtype mismatch & Cast prefix embeddings to model dtype \\
\midrule
\texttt{token\_type\_ids} passed to DistilBERT (rejected) & Filtered \texttt{token\_type\_ids} from input dict \\
\midrule
Dialogue model required personality cache (none existed) & Built 414-NPC cache from training data \\
\midrule
\texttt{src/data/datasets.py} shadowed HF \texttt{datasets} library & Renamed to \texttt{dialogue\_data.py}; updated 7 imports \\
\midrule
\texttt{max\_seq\_len=1024} caused all labels $-100$ (NaN loss) & Increased to 2048 in response/joint configs \\
\midrule
\texttt{best\_dir} missing at save time (\texttt{FileNotFoundError}) & Added \texttt{mkdir(parents=True, exist\_ok=True)} \\
\midrule
\texttt{Gemma4ClippableLinear} layers rejected by PEFT & Changed \texttt{target\_modules} to \texttt{"all-linear"} \\
\midrule
Gemma 4 tokenizer lacks default \texttt{chat\_template} & Set Gemma template with \texttt{<start\_of\_turn>} markers \\
\midrule
Gemma 4 eval OOM on 24\,GB A30 & Set \texttt{PYTORCH\_CUDA\_ALLOC\_CONF=expandable\_segments:True} \\
\midrule
trl 1.3.0 renamed three SFTTrainer parameters & Migrated \texttt{tokenizer}, \texttt{dataset\_text\_field}, \texttt{max\_seq\_length} \\
\bottomrule
\end{tabular}
\end{table*}

\end{document}
